\documentclass[11pt]{article}
\usepackage{amsmath,amssymb,amsthm}
\newtheorem{theorem}{Theorem}
\newtheorem{lemma}{Lemma}
\title{Two Requested Words Already Obstruct a Shared Compact Cauchy Pool}
\author{}
\date{}
\begin{document}
\maketitle

\paragraph{Status and scope.}
This is a complete written argument and a finite exact rational check,
not a Lean theorem. It concerns one fixed positive amplitude pool on
three simple Cauchy poles. It does not permit more pole slots, varying
amplitudes for different requested words, additive query biases, or
query-dependent masks. It does not establish either accepted global
HeadComplexity target.

Fix $0\leq p_1<p_2<p_3\leq a$, where $a>0$, and positive query
times $b\leq t_1<t_2<t_3$, where $b>3a$. For $v_h>0$ and
$\tau\in\{-1,1\}^3$, write
\[
 F_\tau(t)=\sum_{h=1}^3\frac{\tau_hv_h}{t-p_h}.
\]
The amplitude vector $v$ is shared among all choices of $\tau$.

\begin{theorem}[Alternation excludes every one-change word]
If some $F_\tau$ has a strict alternating sign word on
$(t_1,t_2,t_3)$, then the strict words supplied by the entire shared
pool belong to the four-element set consisting of that alternating
word, its negative, and the two constant words. In particular no
shared amplitude vector can realize both an alternating word and a
nonconstant word with exactly one sign change.
\end{theorem}
\begin{proof}
Let $P(t)=\prod_h(t-p_h)$. The numerator $Q=P F_\tau$ has degree at
most two. Strict alternation forces one zero in each of
$(t_1,t_2)$ and $(t_2,t_3)$, since $P$ is positive there.
Consequently $Q(t)=K(t-u)(t-v)$ for two distinct $u,v>b$ and
$K\ne0$. Divide the score by $K$; this changes at most its overall
sign and a common positive amplitude scale. Its residues are
\[
 (r_1,-r_2,r_3),\qquad
 r_1=\frac{(u-p_1)(v-p_1)}{(p_2-p_1)(p_3-p_1)},\quad
 r_2=\frac{(u-p_2)(v-p_2)}{(p_2-p_1)(p_3-p_2)},\quad
 r_3=\frac{(u-p_3)(v-p_3)}{(p_3-p_1)(p_3-p_2)}.
\]
All three displayed magnitudes are positive. We claim that for every
$t\geq b$,
\[
 \frac{r_2}{t-p_2}>\frac{r_1}{t-p_1},\qquad
 \frac{r_2}{t-p_2}>\frac{r_3}{t-p_3}.                 \tag{1}
\]

For the first inequality, consider
\[
 f(p)=\frac{(u-p)(v-p)}{(p_3-p)(t-p)}
 \quad (p_1\leq p\leq p_2).
\]
It is positive, and its logarithmic derivative satisfies
\[
 \frac{f'(p)}{f(p)}
 =-\frac1{u-p}-\frac1{v-p}+\frac1{p_3-p}+\frac1{t-p}
 >\frac1a-\frac2{b-a}>0.
\]
Thus $f(p_2)>f(p_1)$, which is the first inequality in (1) after
division by $p_2-p_1$. For the second one, use
\[
 g(p)=\frac{(u-p)(v-p)}{(p-p_1)(t-p)}
 \quad (p_2\leq p\leq p_3).
\]
Now
\[
 \frac{g'(p)}{g(p)}
 =-\frac1{u-p}-\frac1{v-p}-\frac1{p-p_1}+\frac1{t-p}
 < -\frac1a+\frac1{b-a}<0.
\]
Therefore $g(p_2)>g(p_3)$, giving the other inequality in (1).

There are eight coefficient-sign choices. The two choices
$(+,-,+)$ and $(-,+,-)$ give the original score and its negative.
All-positive and all-negative choices give constant signs. In each
of the remaining four choices, the middle contribution has the
same sign as one outer contribution and, by (1), strictly dominates
the other outer contribution. These four choices therefore also
have constant sign for every $t\geq b$. This proves the theorem.
\end{proof}

\paragraph{Why the shared-amplitude formulation is exact.}
For two requested sign words, one may instead choose an unrestricted
real coefficient vector $c$ for the first word and ask the second
word to use $\operatorname{diag}(S)c$, with $S_h\in\{-1,1\}$.
When all coefficients are nonzero this is equivalent to the shared
positive amplitude formulation, by setting $v_h=|c_h|$. Allowing
zero coefficients cannot repair the alternating three-query word:
two or fewer simple Cauchy terms have numerator degree at most one.

\paragraph{Complete finite test.}
The companion script uses
\[
 C_{ih}=\bigl(1-10^{-4}(i+h+2)\bigr)^{-p},
 \qquad 0\leq i,h\leq2,
\]
for $p=1$ and $p=2$. For every $S\in\{-1,1\}^3$, it enumerates
the strict output words of the stacked matrix
$[C;C\operatorname{diag}(S)]$. The union contains exactly $48$
of the $64$ ordered word pairs for either kernel power. For example,
the first word $(-,+,-)$ and the second word $(+,-,-)$ are absent.
The theorem above proves this absence uniformly for the simple
Cauchy kernel in the stated compact geometry; the squared-kernel
count is only the displayed finite exact result.

Here is why the enumeration is complete even for degenerate
hyperplane arrangements. The stacked matrix $A$ has full column
rank because $C$ is invertible. Any strict word can be rescaled to
make all its signed inequalities at least one. Starting from such
a feasible point, if the active rows have rank below three, move
along a direction annihilating all active rows until a new
independent constraint becomes active. A finite endpoint exists in
at least one direction because $A$ has full column rank. Repeating
gives three independent active rows. The corresponding point solves
$A_Ic=\eta$ for a three-element row set $I$ and
$\eta\in\{-1,1\}^3$. The script enumerates all such nonsingular
systems and directly checks every recorded rational score. It also
checks (1) at $108$ rational configurations, supplementing the
written inequalities rather than replacing them.

\paragraph{Remaining route.}
The obstruction is square and small: three poles and three queries,
with exactly shared magnitudes. A rectangular pool with additional
slots, a carefully controlled alphabet of magnitude changes, or a
different decoder can escape its hypotheses. No asymptotic capacity
claim follows from the finite count or from this theorem alone.
\end{document}
